\section{Some topological concepts in metric spaces}
\label{sec:metric-topology}
Throughout this section, $k$ is a positive integer. In $\R^k$,
$d$ without a subscript denotes Euclidean distance. In an arbitrary
metric space $(S,d)$, write
\[
 B(x,r)=\{y\in S:d(y,x)<r\}\qquad(r>0).
\]
The ambient space matters: complements, interiors, and closures are
taken in the space specified in the exercise. We write $\overline E$
for the book's closure $E^-$, and $E^\circ$ for its interior.

We will use the following elementary ball fact. If $y\in B(x,r)$,
then $\rho=r-d(y,x)>0$. For $z\in B(y,\rho)$, the triangle
inequality gives
\[
 d(z,x)\le d(z,y)+d(y,x)<\rho+d(y,x)=r.
\]
Thus $B(y,\rho)\subseteq B(x,r)$, which proves that every open
ball in every metric space is open. No assertion about the boundary
of a ball in a general metric space is being assumed.

\exercise{13.1}{Two complete metrics on finite-dimensional space}
\goal For $x,y\in\R^k$, verify that
\[
 d_1(x,y)=\max_{1\le j\le k}\abs{x_j-y_j},\qquad
 d_2(x,y)=\sum_{j=1}^k\abs{x_j-y_j}
\]
are metrics, and prove that each metric is complete.
\idea Check the three metric axioms directly. For completeness,
first obtain a limit of each coordinate, then combine the finitely
many coordinate estimates using one cutoff.

\step{(a), first axiom: finite nonnegative distances, with the correct zeros.}
Each $\abs{x_j-y_j}$ is a finite nonnegative real number. A finite
maximum and a finite sum of these quantities are therefore finite
and nonnegative. If $x=y$, every coordinate difference is zero,
and both distances are zero. Conversely, if the maximum is zero,
then every nonnegative coordinate difference is at most zero and
so is zero. If the sum is zero, every nonnegative summand must
also be zero. In either case $x_j=y_j$ for all $j$, hence $x=y$.
Equivalently, if $x\ne y$, at least one difference is positive,
and both distances are positive.

\step{(a), second and third axioms: symmetry and triangle inequalities.}
The identity $\abs{x_j-y_j}=\abs{y_j-x_j}$ gives symmetry for
both formulas. For any $x,y,z\in\R^k$, the real triangle
inequality gives
\[
 \abs{x_j-z_j}\le\abs{x_j-y_j}+\abs{y_j-z_j}.
\]
For $d_1$, the right side is at most $d_1(x,y)+d_1(y,z)$ for
every $j$. Taking the maximum of the left sides proves
$d_1(x,z)\le d_1(x,y)+d_1(y,z)$. For $d_2$, summing the
coordinate inequalities gives
\[
 d_2(x,z)\le\sum_{j=1}^k\abs{x_j-y_j}
              +\sum_{j=1}^k\abs{y_j-z_j}
 =d_2(x,y)+d_2(y,z).
\]
This checks all metric axioms.

\step{(b), coordinate limits: start with a Cauchy sequence in either metric.}
Let $(x^{(n)})$ be Cauchy for $d_i$, where $i=1$ or $i=2$.
For every coordinate $j$,
\[
 \abs{x_j^{(m)}-x_j^{(n)}}\le d_i(x^{(m)},x^{(n)}).
\]
For $d_1$, this is the defining maximum bound; for $d_2$, it is
the bound of one nonnegative summand by the whole sum. Given
$\varepsilon>0$, the Cauchy cutoff for $d_i$ therefore also
makes every coordinate difference smaller than $\varepsilon$.
Each coordinate sequence is Cauchy in $\R$. Completeness of
$\R$ gives a real limit $x_j$ for each $j$. Put
$x=(x_1,\ldots,x_k)\in\R^k$.

\step{(b), convergence in the maximum metric.}
Fix $\varepsilon>0$. For each $j$, choose $N_j$ so that
$n>N_j$ implies $\abs{x_j^{(n)}-x_j}<\varepsilon$.
Take $N=\max\{N_1,\ldots,N_k\}$. If $n>N$, every one of
the finitely many differences is strictly less than $\varepsilon$.
Their maximum is therefore strictly less than $\varepsilon$:
$d_1(x^{(n)},x)<\varepsilon$.

\step{(b), convergence in the sum metric.}
For $d_2$, instead choose $N_j$ to give coordinate errors smaller
than $\varepsilon/k$, which is positive because $k\ge1$.
Beyond the maximum of these cutoffs,
\[
 d_2(x^{(n)},x)=\sum_{j=1}^k\abs{x_j^{(n)}-x_j}
 <k\frac{\varepsilon}{k}=\varepsilon.
\]
Thus a Cauchy sequence for either metric has a limit in that same
metric and in the same space $\R^k$.
\answer{Both $d_1$ and $d_2$ are complete metrics on $\R^k$.}

\exercise{13.2}{Euclidean distance and convergence of coordinates}
\goal Prove both inequalities in Lemma~13.3(1), then prove that
a sequence in $\R^k$ converges if and only if every coordinate
sequence converges in $\R$.

\step{(a) Compare one summand, the sum, and its largest summand.}
For $x,y\in\R^k$, let
$M=\max_{1\le i\le k}\abs{x_i-y_i}$. Every square is
nonnegative, so for each $j$,
\[
 \abs{x_j-y_j}^2\le\sum_{i=1}^k\abs{x_i-y_i}^2\le kM^2.
\]
The first inequality retains one summand and discards only
nonnegative summands. The second replaces each of the $k$
summands by the upper bound $M^2$. Taking nonnegative square
roots, which preserves order, gives
\[
 \boxed{\abs{x_j-y_j}\le d(x,y)\le
 \sqrt{k}\max_{1\le i\le k}\abs{x_i-y_i}.}
\]

\step{(b), forward direction: Euclidean convergence gives coordinate convergence.}
Suppose $x^{(n)}\to x$ in Euclidean distance. Given
$\varepsilon>0$, there is $N$ such that
$d(x^{(n)},x)<\varepsilon$ for $n>N$. The first inequality
in (a) then gives
$\abs{x_j^{(n)}-x_j}<\varepsilon$ for every $j$ and $n>N$.
Thus each $x_j^{(n)}\to x_j$ in $\R$.

\step{(b), reverse direction: select one cutoff for all the coordinates.}
Suppose $x_j^{(n)}\to x_j\in\R$ for $j=1,\ldots,k$,
and define $x=(x_1,\ldots,x_k)$. Given $\varepsilon>0$,
choose $N_j$ such that
\[
 n>N_j\quad\Longrightarrow\quad
 \abs{x_j^{(n)}-x_j}<\frac{\varepsilon}{\sqrt{k}}.
\]
Set $N=\max_jN_j$. For $n>N$, the maximum of all coordinate
errors is smaller than $\varepsilon/\sqrt{k}$. The second
inequality in (a) gives $d(x^{(n)},x)<\varepsilon$.
\answer{The distance bounds are proved, and coordinate convergence
is equivalent to Euclidean convergence.}

\exercise{13.3}{The supremum metric on bounded sequences}
\goal Let $B$ be the set of all bounded real sequences. Prove
that
\[
 d(x,y)=\sup_{j\in\N}\abs{x_j-y_j}
\]
is a metric on $B$, and decide whether the analogous infinite sum
defines a metric on all of $B$.

\step{(a) First verify that the supremum is a finite real number.}
For $x,y\in B$, choose finite nonnegative bounds $C_x,C_y$
with $\abs{x_j}\le C_x$ and $\abs{y_j}\le C_y$ for all $j$.
Then
\[
 0\le\abs{x_j-y_j}\le\abs{x_j}+\abs{y_j}\le C_x+C_y.
\]
The set of coordinate differences is nonempty and bounded above,
so its supremum exists and belongs to $[0,C_x+C_y]$.

\step{(a) Check the zeros, symmetry, and triangle inequality.}
If $x=y$, every difference is zero and so is the supremum.
If $d(x,y)=0$, then each $0\le\abs{x_j-y_j}\le0$, implying
$x_j=y_j$ for all $j$, hence $x=y$. Symmetry follows coordinate
by coordinate from symmetry of absolute value. Finally, for every $j$,
\[
 \abs{x_j-z_j}\le\abs{x_j-y_j}+\abs{y_j-z_j}
 \le d(x,y)+d(y,z).
\]
The final expression is an upper bound for all coordinate
differences between $x$ and $z$. Their least upper bound is thus
at most that expression. This is the triangle inequality for $d$.

\step{(b) A sum can be infinite even when both sequences are bounded.}
Take $x=(0,0,0,\ldots)$ and $y=(1,1,1,\ldots)$.
Both belong to $B$, but
\[
 \sum_{j=1}^{\infty}\abs{x_j-y_j}
 =\sum_{j=1}^{\infty}1=+\infty,
\]
since its $n$th partial sum is $n$. A metric in Definition~13.1
must give a finite real distance between each pair of points.
\answer{The supremum formula is a metric on $B$. The sum formula
is not a metric on all of $B$, because it need not be finite.}

\exercise{13.4}{Unions and finite intersections of open sets}
\goal Prove that arbitrary unions and finite intersections of
open subsets of a metric space $(S,d)$ are open.

\step{(iii) For a point of a union, use one set that contains it.}
Let $U=\bigcup_{\lambda\in\Lambda}U_\lambda$, where each
$U_\lambda$ is open. If $x\in U$, there is a $\lambda_0$ with
$x\in U_{\lambda_0}$. Since that set is open, some $r>0$
satisfies $B(x,r)\subseteq U_{\lambda_0}\subseteq U$.
Thus every point of $U$ is interior, which says $U$ is open.
If the collection is empty, its union is the empty set. That set
is open because there is no point for which the interior condition
could fail.

\step{(iv) For a point of a finite intersection, take the smallest radius.}
Let $V=U_1\cap\cdots\cap U_m$, where $m\ge1$ and all the
$U_i$ are open. For $x\in V$, choose $r_i>0$ with
$B(x,r_i)\subseteq U_i$. The number
$r=\min\{r_1,\ldots,r_m\}$ is positive because this is a finite
nonempty set of positive numbers. If $d(y,x)<r$, then
$d(y,x)<r_i$ for every $i$, so $y\in U_i$ for every $i$.
Hence $B(x,r)\subseteq V$, proving openness.
The intersection of no subsets of $S$ is $S$, which is open:
every ball with center in $S$ is a subset of $S$.

An infinite intersection need not be open: in $\R$,
$\bigcap_{n\ge1}(-1/n,1/n)=\{0\}$, and no positive-radius
interval about $0$ fits in $\{0\}$. This explains why the
intersection statement requires finiteness.
\answer{Arbitrary unions and finite intersections preserve openness.}

\exercise{13.5}{De Morgan's law and intersections of closed sets}
\goal Verify the stated complement identity and deduce that
an arbitrary intersection of closed subsets of $S$ is closed.

\step{(a) Prove equality by testing membership of a point.}
For any $x\in S$, the following statements are equivalent:
\[
 \begin{aligned}
 x\in\bigcap_{U\in\mathcal U}(S\setminus U)
 &\quad\Longleftrightarrow\quad
 x\notin U\text{ for every }U\in\mathcal U\\
 &\quad\Longleftrightarrow\quad
 x\notin\bigcup_{U\in\mathcal U}U\\
 &\quad\Longleftrightarrow\quad
 x\in S\setminus\bigcup_{U\in\mathcal U}U.
 \end{aligned}
\]
Since every point has the same membership on the two sides,
the sets are equal. For the empty collection, the left side is
$S$ and the right side is $S\setminus\varnothing=S$, so the
convention is consistent.

\step{(b) Take complements, then apply the open-union result.}
Let $F_\lambda$ be closed subsets of $S$, and put
$F=\bigcap_{\lambda\in\Lambda}F_\lambda$.
A point fails to belong to this intersection exactly when it
fails to belong to at least one $F_\lambda$. Thus
\[
 S\setminus F
 =\bigcup_{\lambda\in\Lambda}(S\setminus F_\lambda).
\]
Each complement on the right is open by the definition of a
closed set. Their union is open by Exercise~13.4, so $F$ is
closed by that same definition. If $\Lambda$ is empty, $F=S$
and its complement is empty and open.
\answer{The complement identity holds, and every intersection of
closed sets is closed.}

\exercise{13.6}{Four equivalent descriptions involving closure}
\goal Prove all four assertions of Proposition~13.9 for a subset
$E$ of an arbitrary metric space $(S,d)$.
\idea First derive a ball characterization directly from the
definition of closure as an intersection of closed supersets.
Then obtain the sequence characterizations without assuming them.

\step{Preparation: closure is a closed superset, and a ball detects its points.}
The family of closed sets containing $E$ is nonempty, because
$S$ belongs to it. Every set in that family contains $E$;
therefore their intersection $\overline E$ contains $E$.
Exercise~13.5 shows that the intersection is closed.

For any $x\in S$, we claim
\[
 x\in\overline E
 \quad\Longleftrightarrow\quad
 B(x,r)\cap E\ne\varnothing\text{ for every }r>0.\tag{*}
\]
If some ball $B(x,r)$ misses $E$, then its complement is a
closed set containing $E$, since balls are open. This closed
superset excludes $x$, so $x\notin\overline E$.
Conversely, if $x\notin\overline E$, then at least one closed
set $C$ containing $E$ excludes $x$; otherwise $x$ would belong
to their intersection. The open set $S\setminus C$ contains
$x$, so it contains some $B(x,r)$, and that ball misses $E$.
These two implications prove (*).

\step{(a) A set is closed exactly when it equals its closure.}
If $E$ is closed, it is one of the closed supersets in the
defining intersection, so $\overline E\subseteq E$.
The reverse inclusion was already proved, giving equality.
If $E=\overline E$, then $E$ is closed because $\overline E$
was proved closed. This proves both directions.

\step{(c) Closure points are exactly limits of sequences of points of $E$.}
We establish (c) before (b) so that there is no circular argument.
If $x\in\overline E$, apply (*) with radius $1/n$ to choose
$x_n\in E$ with $d(x_n,x)<1/n$. For any $\varepsilon>0$,
choose an integer $N>1/\varepsilon$. If $n>N$, then
$d(x_n,x)<1/n<\varepsilon$. Hence $x_n\to x$.

Conversely, if $x_n\in E$ and $x_n\to x\in S$, fix $r>0$.
Convergence gives some index $n$ with $d(x_n,x)<r$, so
$x_n\in B(x,r)\cap E$. Every ball about $x$ thus meets $E$,
and (*) gives $x\in\overline E$. For $E=\varnothing$, no such
sequence exists and its closure is empty, so the equivalence
still holds.

\step{(b) Closed sets are exactly those containing limits of their convergent sequences.}
If $E$ is closed and $x_n\in E$ tends to a point $x\in S$,
part (c) gives $x\in\overline E$, and part (a) gives
$\overline E=E$, so $x\in E$.
Conversely, suppose $E$ contains every such limit. For every
$x\in\overline E$, part (c) supplies a sequence in $E$
converging to $x$, so the assumption gives $x\in E$.
Thus $\overline E\subseteq E$. Combined with the always-valid
inclusion $E\subseteq\overline E$, this gives equality, and
part (a) makes $E$ closed.

\step{(d) Boundary points see both the set and its complement.}
By definition, the boundary is $\overline E\setminus E^\circ$.
A point $x$ lies outside $E^\circ$ exactly when no ball about
$x$ is contained in $E$. Indeed, the existence of such a ball
would also put $x$ in $E$, since its center belongs to the ball,
and would make $x$ interior. Saying no ball is contained in $E$
is equivalent to saying every ball meets $S\setminus E$.
By (*) applied to this complement,
\[
 x\notin E^\circ\quad\Longleftrightarrow\quad
 x\in\overline{S\setminus E}.
\]
Consequently
\[
 \overline E\setminus E^\circ
 =\overline E\cap\overline{S\setminus E},
\]
which is precisely the assertion in (d).
\answer{All four claims follow from the closure definition and
the ball criterion (*), with no sequence characterization assumed
in advance.}

\exercise{13.7}{Decomposing an open real set into disjoint intervals}
\goal Express every open $U\subseteq\R$ as a finite or countably
infinite disjoint union of open intervals. Unbounded intervals are
allowed; the empty set is the union of an empty family.

\step{1. Group points that have no gap of $U$ between them.}
For $x,y\in U$, write $x\sim y$ when the whole closed segment
with endpoints $x,y$ lies in $U$. This relation is reflexive,
because $[x,x]=\{x\}\subseteq U$, and symmetric because
the segment is unchanged when its endpoints are interchanged.
For transitivity, if $x\sim y$ and $y\sim z$, every point
between $x$ and $z$ lies in the union of the segments joining
$x$ to $y$ and $y$ to $z$. That union is contained in $U$.
Thus $x\sim z$, so $\sim$ is an equivalence relation.
Its classes partition $U$ into nonempty pairwise disjoint sets.

\step{2. Show that each class is an interval and is open.}
Let $I$ be a class, and let $p<q$ belong to it. The equivalence
relation gives $[p,q]\subseteq U$. For $p<t<q$, the segment
$[p,t]$ is contained in $[p,q]$, so $p\sim t$ and $t\in I$.
Thus every point between two points of $I$ belongs to $I$.

For $z\in I$, openness of $U$ gives $r>0$ with
$(z-r,z+r)\subseteq U$. Every $t$ in that interval has its
segment to $z$ inside the interval, hence inside $U$. Therefore
$t\sim z$ and $t\in I$. We have proved
$(z-r,z+r)\subseteq I$, so $I$ is itself open.

\step{3. Identify the endpoints, including infinite endpoints.}
Let $a=\inf I$ and $b=\sup I$ in the extended real sense;
use $a=-\infty$ or $b=+\infty$ if the corresponding bound
does not exist in $\R$. Since $I$ is nonempty and open,
it contains an interval around one of its points, and $a<b$.
For $a<t<b$, the definitions of infimum and supremum give
$p,q\in I$ with $p<t<q$; otherwise $t$ would be a lower or
upper bound contradicting $a<t<b$. The interval property puts
$t\in I$. Conversely, a finite endpoint cannot lie in $I$:
openness there would supply a point of $I$ still smaller than
$a$ or still larger than $b$. Hence $I=(a,b)$, with the
usual interpretations $(a,+\infty)$, $(-\infty,b)$, or $\R$.

\step{4. Prove that there are at most countably many classes.}
Fix an enumeration $q_1,q_2,\ldots$ of $\Q$. Each nonempty
open interval $I$ contains a rational by density. Assign to $I$
the least index $m(I)$ for which $q_{m(I)}\in I$.
This least index exists by the least-element property of $\N$.
Distinct classes have distinct assigned indices, because they
are disjoint and cannot contain the same rational.
Thus the classes inject into $\N$. If there are infinitely
many, list their assigned indices in increasing order and use
that order to list the classes; every assigned index appears,
since only finitely many positive integers precede it.

\answer{The equivalence classes give the required disjoint interval
decomposition. The family may be empty, finite, or countably
infinite, and intervals may be unbounded.}

\exercise{13.8}{Intervals, balls, half-spaces, and examples that are neither}
\goal Verify the assertions in the first paragraph of Example~3
and all the assertions in Example~4. In the interval assertions
assume $a<b$; in the ball assertions assume $r>0$ and $k\ge1$.
We address the degenerate cases separately below.

\step{(a) Open and closed intervals in the real line.}
For $x\in(a,b)$, let $\rho=\min\{x-a,b-x\}>0$.
If $\abs{y-x}<\rho$, then $a<y<b$. Thus $(a,b)$ is open.
The complement of $[a,b]$ is
$(-\infty,a)\cup(b,+\infty)$. Each ray is open: for a point
strictly on its allowed side, half its distance to the endpoint
provides a ball staying on that side. Their union is open,
so $[a,b]$ is closed.

Every $x\in(a,b)$ is interior to $[a,b]$ by the first estimate.
The endpoint $a$ is not interior: every ball about it contains
$a-\rho/2<a$ for its radius $\rho>0$. Similarly $b$ is not
interior because $b+\rho/2$ is outside $[a,b]$.
Hence $[a,b]^\circ=(a,b)$.

The closed set $[a,b]$ contains $(a,b)$, so
$\overline{(a,b)}\subseteq[a,b]$. Every ball about $a$ meets
$(a,b)$: choose a displacement
$0<t<\min\{\rho,b-a\}$ and use $a+t$. The point $b-t$
does the same for $b$. The ball criterion of Exercise~13.6
therefore puts both endpoints in the closure. Interior points
already belong to the set. Thus
$\overline{(a,b)}=[a,b]=\overline{[a,b]}$.
Subtracting the interior gives boundary $\{a,b\}$ for both
intervals.

If $a=b$, the open interval is empty and has empty closure and
boundary. The closed interval is the singleton $\{a\}$, which
is closed, has empty interior, and has boundary $\{a\}$.
This is why the boundary assertion for both intervals presumes
$a<b$.

\step{(b), balls: openness, closedness, and the possible boundary points.}
For a center $c\in\R^k$, put
\[
 D=\{x:d(x,c)<r\},\qquad K=\{x:d(x,c)\le r\}.
\]
The opening ball argument of this section shows $D$ is open.
If $x\notin K$, let $\rho=d(x,c)-r>0$. For
$d(y,x)<\rho$, the triangle inequality rearranges to
\[
 d(y,c)\ge d(x,c)-d(x,y)>d(x,c)-\rho=r.
\]
Thus a ball about $x$ stays outside $K$, so the complement
of $K$ is open and $K$ is closed.

If $d(x,c)<r$, then $x$ is interior to both sets, using the
radius $r-d(x,c)$. If $d(x,c)>r$, the preceding complement
ball shows $x$ lies outside the closure of either set.
Only points with $d(x,c)=r$ can therefore be boundary points.

\step{(b), spheres: every point on the sphere really is a boundary point.}
Let $d(x,c)=r>0$, and let $\varepsilon>0$. Choose
$0<t<\min\{1,\varepsilon/r\}$. Define
\[
 x_-=c+(1-t)(x-c),\qquad x_+=c+(1+t)(x-c).
\]
The Euclidean distance formula gives
\[
 d(x_-,c)=(1-t)r<r,\quad d(x_+,c)=(1+t)r>r,
 \quad d(x_-,x)=d(x_+,x)=tr<\varepsilon.
\]
Thus every ball about $x$ meets $D$ (at $x_-$) and the
complement of $K$ (at $x_+$). It meets both either set and
its complement, so Exercise~13.6(d) puts $x$ on both boundaries.
Consequently both boundaries are exactly $\{x:d(x,c)=r\}$.
This radial argument uses Euclidean space; the same sphere
assertion is not valid for all metrics. If $r=0$, the open ball
is empty with empty boundary, while the closed ball is the
singleton $\{c\}$ with boundary $\{c\}$.

\step{(b), strict coordinate inequalities give open half-spaces and quadrants.}
In $\R^2$, consider $H=\{x:x_1>0\}$ and
$Q=\{x:x_1>0,\ x_2>0\}$. For $x\in H$, choose
$\rho=x_1/2>0$. Since $\abs{y_1-x_1}\le d(y,x)$,
every $y$ in $B(x,\rho)$ satisfies
$y_1>x_1-\rho=x_1/2>0$. Thus $H$ is open.
For $x\in Q$, choose $\rho=\min\{x_1,x_2\}/2$.
The same coordinate estimate, separately for $j=1,2$,
gives $y_j>x_j-\rho\ge x_j/2>0$. Hence $Q$ is open.

For $H_0=\{x:x_1\ge0\}$, its complement is
$\{x:x_1<0\}$, an open half-space by the same radius argument
with $-x_1$ in place of $x_1$. Thus $H_0$ is closed.
For $Q_0=\{x:x_1\ge0,\ x_2\ge0\}$, the complement is
$\{x:x_1<0\}\cup\{x:x_2<0\}$, a union of open half-spaces.
Thus $Q_0$ is closed as asserted.

\step{(b), the half-open interval is neither open nor closed.}
The point $0\in[0,1)$ is not interior, since every ball about
$0$ contains negative points. Hence $[0,1)$ is not open.
The sequence $1-1/n$ belongs to $[0,1)$ and tends to $1$,
which is outside the set. The closed-set sequence criterion
of Exercise~13.6(b) shows it is not closed.

\step{(b), the mixed quadrant is neither open nor closed.}
Let $M=\{x:x_1>0,\ x_2\ge0\}$. The point $(1,0)$ belongs
to $M$, but every ball of radius $\rho>0$ about it contains
$(1,-\rho/2)$, outside $M$. Thus $M$ is not open.
The sequence $(1/n,0)$ lies in $M$ and tends in Euclidean
distance to $(0,0)$, outside $M$. It follows that $M$ is not
closed either.
\answer{All the interval, Euclidean ball, half-space, quadrant,
and mixed examples have the stated properties, with the positive
radius and nondegenerate interval conditions made explicit.}

\exercise{13.9}{Three closures in the real line}
\goal Find the closures of
$\{1/n:n\in\N\}$, $\Q$, and $\{q\in\Q:q^2<2\}$ in $\R$.

\step{(a) Include the reciprocal terms and their one additional limit.}
Let $A=\{1/n:n\in\N\}$. Every point of $A$ belongs to its
closure, and $1/n\to0$ gives $0\in\overline A$ by
Exercise~13.6(c). We show no other point belongs.
If $x<0$, the ball of radius $-x/2$ about $x$ is negative
and misses $A$. If $x>1$, the radius $(x-1)/2$ gives a
ball wholly above $1$, which also misses $A$.

It remains to exclude $0<x\le1$ with $x\notin A$.
Choose $N$ so large that $1/N<x/2$. All terms $1/n$ for
$n\ge N$ are then below $x/2$. The finite set
$\{1,1/2,\ldots,1/(N-1)\}$ has no member equal to $x$.
Its finitely many positive distances from $x$ have a positive
minimum $\eta$. Here $N\ge2$ because $1/N<x/2\le1/2$.
A ball about $x$ of radius $\min\{x/2,\eta\}/2$
misses the finite set and all the remaining reciprocals.
The ball criterion excludes $x$. Therefore
\[
 \boxed{\overline A=A\cup\{0\}.}
\]

\step{(b) Density fills the whole real line.}
For any $x\in\R$ and $r>0$, the interval $(x-r,x+r)$
contains a rational by density of $\Q$. Hence every real $x$
satisfies the closure ball criterion, giving
$\boxed{\overline\Q=\R}$.

\step{(c) The rational points between the two square-root endpoints.}
For real $q$, the inequality $q^2<2$ is equivalent to
$-\sqrt2<q<\sqrt2$. Thus the given set is
$C=\Q\cap(-\sqrt2,\sqrt2)$.
If $-\sqrt2<x<\sqrt2$, then every neighborhood of $x$
contains a smaller nonempty interval lying inside
$(-\sqrt2,\sqrt2)$; density supplies a rational in that interval.
For $x=\sqrt2$, an interval
$(\sqrt2-t,\sqrt2)$ with $0<t<\min\{r,2\sqrt2\}$
lies inside the radius-$r$ neighborhood and inside the allowed
range, and contains a rational. The interval
$(-\sqrt2,-\sqrt2+t)$ proves the same statement for $x=-\sqrt2$.
Therefore the whole closed interval belongs to $\overline C$.

If $x>\sqrt2$, take radius $(x-\sqrt2)/2$; if
$x<-\sqrt2$, take radius $(-\sqrt2-x)/2$. The resulting
ball misses the closed interval and hence misses $C$.
No outside point belongs to the closure, so
\[
 \boxed{\overline C=[-\sqrt2,\sqrt2].}
\]
\answer{The closures are respectively
$\{0\}\cup\{1/n:n\in\N\}$, $\R$, and $[-\sqrt2,\sqrt2]$.}

\exercise{13.10}{Three sets with no interior points}
\goal Prove that the reciprocal set, the rationals, and the
Cantor set all have empty interior in $\R$.

\step{(a) A rational-only set cannot contain a real interval.}
Every reciprocal $1/n$ is rational. Every open interval contains
an irrational number. For completeness, to find one in $(a,b)$,
choose a rational $q$ in $(a-\sqrt2,b-\sqrt2)$ by rational
density; then $q+\sqrt2\in(a,b)$ is irrational, since its
being rational would make $\sqrt2=(q+\sqrt2)-q$ rational.
Thus every ball about any point of the reciprocal set contains
a point outside that set. No point of the set is interior.

\step{(b) Apply the same irrational-density argument to $\Q$.}
For any $q\in\Q$ and $r>0$, the preceding argument gives an
irrational in $(q-r,q+r)$. That ball is not a subset of $\Q$.
Thus $\Q$ also has no interior points. Notice that dense closure
does not imply nonempty interior: Exercise~13.9 gave
$\overline\Q=\R$.

\step{(c) The intervals remaining at late Cantor stages are too short.}
Use the numbering in Example~5: $F_1=[0,1]$, and $F_n$ is
the union of $2^{n-1}$ disjoint closed intervals, each of length
$3^{-(n-1)}$. The Cantor set is $F=\bigcap_{n\ge1}F_n$.
Suppose, for contradiction, that $x\in F^\circ$. Then some
real interval $(x-r,x+r)$, with $r>0$, lies in $F$.
The closed interval $J=[x-r/2,x+r/2]$ has length $r$ and
also lies in $F$, hence in every $F_n$.

Choose $n$ so large that $3^{-(n-1)}<r$. The interval $J$
would have to lie in one of the component intervals of $F_n$.
Indeed, if it met two distinct components, it would also contain
the points of the nonempty removed gap between those components,
because it contains every point between any two of its points.
That would contradict $J\subseteq F_n$.
But an interval of length $r$ cannot be contained in an interval
of smaller length $3^{-(n-1)}$. This is a contradiction.
\answer{All three interiors are empty. The Cantor proof uses the
explicit shrinking intervals of its construction.}

\exercise{13.11}{Compactness and subsequences in Euclidean space}
\goal Show that $E\subseteq\R^k$ is compact if and only if
every sequence of points of $E$ has a subsequence converging to
a point of $E$.
\idea Use Heine--Borel and Bolzano--Weierstrass. For the converse,
show that the stated subsequence property forces both boundedness
and closedness.

\step{Preparation: uniqueness and inheritance of metric limits.}
If $x_n\to x$ and $x_n\to y$, then
$d(x,y)\le d(x,x_n)+d(x_n,y)\to0$, so $d(x,y)=0$ and $x=y$.
Every subsequence of a convergent sequence has the same limit:
strictly increasing selection indices $n_j$ satisfy $n_j\ge j$,
so every eventual distance estimate for the original sequence
also holds for all sufficiently late selected terms.
Also, a convergent sequence in $\R^k$ is bounded: its tail
lies within distance $1$ of its limit $x$, and for each
coordinate the triangle inequality gives
$\abs{x_j^{(n)}}\le\abs{x_j}+1$ on that tail. A finite
maximum accommodates the omitted initial coordinates.

\step{Forward direction: combine boundedness, subsequences, and closedness.}
If $E$ is compact, Heine--Borel (Theorem~13.12) says it is
closed and bounded. Any sequence in $E$ is then bounded.
Bolzano--Weierstrass (Theorem~13.5) gives a subsequence tending
to some point of $\R^k$. Since all its terms lie in the closed
set $E$, Exercise~13.6(b) puts that limit in $E$.

\step{Reverse direction, boundedness: unbounded points would escape every subsequence.}
Assume the stated subsequence property. If $E$ were unbounded,
then for each $n\in\N$ we could choose $x^{(n)}\in E$ with
\[
 \max_{1\le j\le k}\abs{x_j^{(n)}}>n.
\]
For any subsequence $x^{(n_m)}$, the selected indices satisfy
$n_m\ge m$, so its coordinate maximum exceeds $m$.
It is therefore unbounded: for any proposed finite bound $C$,
take $m>C$. It cannot converge, since a convergent sequence
was just shown to be bounded. This contradicts the assumed
property. Thus $E$ is bounded.

\step{Reverse direction, closedness: recover every closure point.}
Let $x\in\overline E$. By Exercise~13.6(c), choose a sequence
$x_n\in E$ converging to $x$. The assumed property supplies
a subsequence $x_{n_m}\to y$ for some $y\in E$.
Inheritance gives $x_{n_m}\to x$ as well, and uniqueness of
the limit gives $x=y\in E$. Thus $\overline E\subseteq E$.
The reverse inclusion always holds, so Exercise~13.6(a)
makes $E$ closed. Having proved both closedness and boundedness,
Heine--Borel gives compactness.

If $E=\varnothing$, it is compact because the empty subfamily
is a finite subcover of every cover of it. There is no sequence
of points of $E$, so the sequential condition is vacuously true.
\answer{The two conditions are equivalent for all subsets of
$\R^k$, including the empty set.}

\exercise{13.12}{Closed subsets and finite unions of compact sets}
\goal In any metric space $(S,d)$, prove that a closed subset
of a compact set is compact, and that a finite union of compact
sets is compact.
\idea Work directly with open covers, so neither result depends
on the finite-dimensional Heine--Borel theorem.

\step{(a) Add one open set to cover the rest of the larger compact set.}
Suppose $E\subseteq F$, $E$ is closed in $S$, and $F$ is
compact. Let $\mathcal U$ be any family of open subsets of
$S$ covering $E$. The set $V=S\setminus E$ is open, and
$\mathcal U\cup\{V\}$ covers $F$: a point of $F$ either
belongs to $E$ and is covered by $\mathcal U$, or lies in $V$.
Compactness of $F$ gives a finite subfamily covering $F$.
Remove $V$ if it was selected. The remaining sets still cover
$E$, since no point of $E$ belongs to $V$. They are a finite
subfamily of $\mathcal U$, so $E$ is compact.

If ``closed subset of $F$'' is interpreted as closed relative
to $F$, the same conclusion holds: relative closedness means
$F\setminus E=F\cap V$ for some open $V\subseteq S$.
Use that $V$ in the preceding cover argument; it still misses
$E$ and covers every point of $F\setminus E$.

\step{(b) Take a finite subcover for each compact piece.}
Let $K=K_1\cup\cdots\cup K_m$ with each $K_i$ compact,
and let $\mathcal U$ be an open cover of $K$.
It also covers each $K_i$. For each $i$, select a finite
subfamily $\mathcal U_i\subseteq\mathcal U$ covering $K_i$.
The family $\mathcal U_1\cup\cdots\cup\mathcal U_m$ is
finite, since a finite union of finite sets is finite, and
it covers every piece, hence their union $K$. This proves
compactness. When $m=0$, the union is empty and is compact
by the empty-subcover argument.
\answer{Both permanence properties hold in every metric space.}

\exercise{13.13}{A compact real set contains its extreme bounds}
\goal For a nonempty compact $E\subseteq\R$, show
$\sup E\in E$ and $\inf E\in E$.

\step{1. Ensure the bounds are finite and use the approximation property.}
Heine--Borel makes $E$ closed and bounded. Since it is also
nonempty, completeness of $\R$ supplies finite real numbers
$M=\sup E$ and $m=\inf E$.
For every $n\in\N$, $M-1/n$ is strictly below the least
upper bound, so it cannot be an upper bound. Choose $u_n\in E$
with
\[
 M-\frac1n<u_n\le M.
\]
Similarly, $m+1/n$ is strictly above the greatest lower bound,
so it cannot be a lower bound. Choose $v_n\in E$ with
\[
 m\le v_n<m+\frac1n.
\]

\step{2. Check convergence, then use closedness.}
The first inequalities give $\abs{u_n-M}<1/n$, so $u_n\to M$.
The second give $\abs{v_n-m}<1/n$, so $v_n\to m$.
Both sequences lie in $E$. Closedness and Exercise~13.6(b)
therefore put both limits in $E$.
\answer{The supremum and infimum are attained: $E$ has a maximum
and a minimum, namely $M$ and $m$.}

\exercise{13.14}{The diameter of a compact set is attained}
\goal For nonempty compact $E\subseteq\R^k$, put
\[
 \delta=\sup\{d(x,y):x,y\in E\}.
\]
Find $x^0,y^0\in E$ such that $d(x^0,y^0)=\delta$.
\idea Approximate the supremum by pairs, select subsequences
for both entries of the pairs, and justify passage to the limit
using the triangle inequality.

\step{1. Verify that the distance supremum is a finite real number.}
Heine--Borel makes $E$ bounded. Choose $C\ge0$ so that every
coordinate of every $x\in E$ has absolute value at most $C$.
For $x,y\in E$, each $\abs{x_j-y_j}\le2C$, so
$d(x,y)\le2C\sqrt{k}$ by Exercise~13.2(a).
The set of distances is nonempty because $E$ is nonempty,
and contains $d(x,x)=0$. Its supremum therefore exists with
$0\le\delta\le2C\sqrt{k}$.

\step{2. Choose pairs whose distances approach the supremum.}
For each $n\in\N$, the real number $\delta-1/n$ is not
an upper bound of the distance set. Hence choose $x_n,y_n\in E$
such that
\[
 \delta-\frac1n<d(x_n,y_n)\le\delta.
\]
This works even when $\delta=0$; any pair then has distance
$0>\delta-1/n$. The displayed error bound gives
$d(x_n,y_n)\to\delta$.

\step{3. Pass to a common subsequence for the two endpoints.}
By Exercise~13.11, some subsequence $x_{n_j}$ tends to a
point $x^0\in E$. Apply the same result to the sequence
$y_{n_j}$ in $E$ to find further indices $j_\ell$ with
$y_{n_{j_\ell}}\to y^0\in E$.
The matching $x_{n_{j_\ell}}$ still tends to $x^0$, since
a subsequence inherits a limit. This further selection is
also a subsequence of the original sequence of pairs because
the composition of two increasing index selections is increasing.

\step{4. Prove that distances pass to the limit.}
Two uses of the triangle inequality give
\[
 \begin{split}
 d(x,y)&\le d(x,x^0)+d(x^0,y^0)+d(y^0,y),\\
 d(x^0,y^0)&\le d(x^0,x)+d(x,y)+d(y,y^0).
 \end{split}
\]
Subtracting $d(x^0,y^0)$ in the first and $d(x,y)$ in the
second combines into
\[
 \abs{d(x,y)-d(x^0,y^0)}\le d(x,x^0)+d(y,y^0).
\]
At the selected pairs, the right side tends to zero. Thus
$d(x_{n_{j_\ell}},y_{n_{j_\ell}})\to d(x^0,y^0)$.
The same distances tend to $\delta$, being a subsequence of
the distance sequence in Step~2. Uniqueness of real limits
gives $d(x^0,y^0)=\delta$.
\answer{Two points of $E$ realize its diameter, including the
zero-diameter case.}

\exercise{13.15}{A closed bounded set that is not compact}
\goal In the bounded-sequence space $(B,d)$ of Exercise~13.3,
let
\[
 F=\left\{x\in B:\sup_{j\in\N}\abs{x_j}\le1\right\}.
\]
Show that $F$ is closed and bounded, but not compact.
\idea Verify closedness using an open complement. Then use
the prescribed cover by open balls of radius $1$, together
with infinitely many points whose pairwise distance is $2$.

\step{(a), boundedness: use the zero sequence as the center.}
The sequence $0=(0,0,\ldots)$ belongs to $B$. For $x\in F$,
\[
 d(x,0)=\sup_j\abs{x_j}\le1.
\]
Thus $F$ lies in the closed metric ball with center $0$
and radius $1$, which is exactly the boundedness condition
stated in the exercise.

\step{(a), closedness: exhibit a ball about each point outside $F$.}
If $x\in B\setminus F$, then $\sup_j\abs{x_j}>1$.
At least one coordinate, say $j_0$, satisfies $\abs{x_{j_0}}>1$;
otherwise $1$ would be an upper bound for all these absolute
values and hence for their supremum. Set
$\rho=(\abs{x_{j_0}}-1)/2>0$.
If $d(x,y)<\rho$, the real triangle inequality gives
\[
 \abs{y_{j_0}}\ge\abs{x_{j_0}}-\abs{x_{j_0}-y_{j_0}}
 >\abs{x_{j_0}}-\rho
 =\frac{\abs{x_{j_0}}+1}{2}>1.
\]
Such a $y$ lies outside $F$. Hence a ball about each point
of the complement stays in the complement, proving it open
and $F$ closed in $B$.

\step{(b), the open cover: give every point its radius-one ball.}
For $x\in F$, define
$U(x)=\{y\in B:d(y,x)<1\}$. These sets are open in $B$
by the general ball fact at the start of this section.
They cover $F$, since $d(x,x)=0<1$ puts each $x\in F$
in its own $U(x)$.

\step{(b), separated points: compute the distance exactly.}
For each $n\in\N$, let $x^{(n)}$ have coordinate $-1$
at index $n$ and coordinate $1$ at every other index.
Every one of its coordinates has absolute value $1$,
so $x^{(n)}\in F$.
If $m\ne n$, the two sequences differ by $2$ in absolute
value at coordinates $m$ and $n$, and agree at all other
coordinates. Therefore
\[
 d(x^{(m)},x^{(n)})=2.
\]
No ball $U(x)$ can contain two distinct such points: if it
contained both, then the metric triangle inequality would give
\[
 2=d(x^{(m)},x^{(n)})
 \le d(x^{(m)},x)+d(x,x^{(n)})<1+1=2,
\]
an impossibility. The strict inequalities in the open-ball
definition are essential here.

\step{(b), exclude every finite subcover.}
A union of $q$ members of this cover can contain at most
$q$ of the separated points, since each member contains at
most one. But $F$ contains the infinitely many distinct
points $x^{(1)},x^{(2)},\ldots$. For example, among its first
$q+1$ separated points at least one must be missed. Thus
no finite subfamily covers $F$.
\answer{$F$ is closed and bounded in the specified metric space,
but this explicit open cover has no finite subcover, so $F$
is not compact. Heine--Borel's closed-and-bounded criterion
is specific to Euclidean space and does not apply to all
metric spaces.}
